%document
\documentclass[10pt]{beamer}
%theme
\usetheme{CambridgeUS}
% packages
\usepackage{color}
\usepackage{listings}
\usepackage[ngerman]{babel}
\usepackage[utf8]{inputenc}
\usepackage{multicol}


% color definitions
\definecolor{mygreen}{rgb}{0,0.6,0}
\definecolor{mygray}{rgb}{0.5,0.5,0.5}
\definecolor{mymauve}{rgb}{0.58,0,0.82}

\lstset{
    backgroundcolor=\color{white},
    % choose the background color;
    % you must add \usepackage{color} or \usepackage{xcolor}
    basicstyle=\footnotesize\ttfamily,
    % the size of the fonts that are used for the code
    breakatwhitespace=false,
    % sets if automatic breaks should only happen at whitespace
    breaklines=true,                 % sets automatic line breaking
    captionpos=b,                    % sets the caption-position to bottom
    commentstyle=\color{mygreen},    % comment style
    % deletekeywords={...},
    % if you want to delete keywords from the given language
    extendedchars=true,
    % lets you use non-ASCII characters;
    % for 8-bits encodings only, does not work with UTF-8
    frame=single,                    % adds a frame around the code
    keepspaces=true,
    % keeps spaces in text,
    % useful for keeping indentation of code
    % (possibly needs columns=flexible)
    keywordstyle=\color{blue},       % keyword style
    % morekeywords={*,...},
    % if you want to add more keywords to the set
    numbers=left,
    % where to put the line-numbers; possible values are (none, left, right)
    numbersep=5pt,
    % how far the line-numbers are from the code
    numberstyle=\tiny\color{mygray},
    % the style that is used for the line-numbers
    rulecolor=\color{black},
    % if not set, the frame-color may be changed on line-breaks
    % within not-black text (e.g. comments (green here))
    stepnumber=1,
    % the step between two line-numbers.
    % If it's 1, each line will be numbered
    stringstyle=\color{mymauve},     % string literal style
    tabsize=4,                       % sets default tabsize to 4 spaces
    % show the filename of files included with \lstinputlisting;
    % also try caption instead of title
    language = Java,
	showspaces = false,
	showtabs = false,
	showstringspaces = false,
	escapechar = ,
}

\def\ContinueLineNumber{\lstset{firstnumber=last}}
\def\StartLineAt#1{\lstset{firstnumber=#1}}
\let\numberLineAt\StartLineAt



\newcommand{\codeline}[1]{
	\alert{\texttt{#1}}
}

% This Document contains the information about this course.

% Authors of the slides
\author{Vincent Gerber, Tilman Hinnerichs}

% Name of the Course
\institute{Java Kurs}

% Fancy Logo
\titlegraphic{\hfill\includegraphics[height=1.25cm]{../templates/fsr_logo_cropped}}


\title{Java}
\subtitle{Introduction to OOP and inheritance}
\date{\today}

\begin{document}

\section{OOP in Java}

\begin{frame}{}
	\begin{center}
		{\huge Object Oriented Programming}
	\end{center}
\end{frame}

\subsection{General information}

\begin{frame}[fragile]{What is OOP?}
	We have a lot of objects around us\\
	\begin{itemize}
		\item They have a state and a behaviour
		\item All have a building plan
	\end{itemize}
	
	Classes $ \rightarrow $ building plan\\
	Objects $ \rightarrow $ entity in real world build from building plan
	
\end{frame}

\begin{frame}[fragile]{How to build a class}
	\begin{lstlisting}
		public class MyClass {
			//Constructor to create the Object
			public MyClass(){
			...
			}
			
			//Attributes to describe state
			private int myAttribute;
			...
			
			//Methods to describe behaviour
			public MyMethod(int myParameter, ...) {
				this.myAttribute = myParameter;
				...
			}
		}
	\end{lstlisting}
\end{frame}

\begin{frame}[fragile]{How to build a student?}
	Attributes describing a student? \\
	\pause 
	\begin{itemize}
		\item Name
		\item Matriculation number
		\item Age
	\end{itemize}	
	\pause
	Methods describing its behaviour?
	\pause
	\begin{itemize}
		\item Change name
		\item Increase age
		\item Calculate 1+1
	\end{itemize}
\end{frame}

\begin{frame}[fragile]{Class Student}
	\begin{lstlisting}
	public class Student {
	//Constructor
	public Student (String name, int matriculationNumber) {
	this.name = name;
	this.matriculationNumber = matriculationNumber;
	}
	
	// Attributes
	private String name; 
	private int matriculationNumber; 
	
	
	// Methods
	public void setName(String name) {
	this.name = name;
	}
	
	public int getMatriculationNumber() {
	return matriculationNumber;
	}
	
	}
	\end{lstlisting}
	
	% What is visible here:
	% Attributes store the state of the object
	% Methods implement the behaviour of the object
	
\end{frame}

\begin{frame}[fragile]{Creation}
	We learned how to declare and assign a primitive datatype.}

\begin{lstlisting}
int a; // declare a
a = 273; // assign 273 to a
\end{lstlisting} 

The creation of an object works similar.

\begin{lstlisting}
Student example = new Student(); 
// create an instance of Student
\end{lstlisting}
The \textbf{object} derived from a \textbf{class} is also called \textbf{instance}.
The variable is called the \textbf{reference}.
\end{frame}

\subsection{Methods}
\begin{frame}[fragile]{Calling a Method}
\begin{lstlisting}
public class Student {

private String name;

public String getName() {
return name;
}
public void setName(String newName) {
name = newName;
}
public void printName(){
System.out.println(this.name);
}
}
\end{lstlisting}
The class \emph{Student} has three methods: \emph{void printName()}.
\end{frame}

\begin{frame}[fragile]{Calling a Method}
\begin{lstlisting}
public class Main {

public static void main(String[] args) {
Student example = new Student(); // creation
example.setName("Jane"); // method call
String name = example.getName(); 
System.out.println(name); // Prints "Jane"

example.printName();
}

}
\end{lstlisting}
You can call a method of an object after its creation with \textbf{reference.methodName();}.
\end{frame}

\begin{frame}[fragile]{Calling a Method}
\begin{lstlisting}
public class Student {

private String name;

public void setName(String newName) {
name = newName;
printName();   // Call own method
this.printName(); // Or this way
}

public void printName() {
System.out.println(name);
}

}
\end{lstlisting}
You can call a method of the own object by simply writing \textbf{methodName();} or \textbf{this.methodName();}
\end{frame}

\begin{frame}[fragile]{Methods with Arguments}

\begin{lstlisting}
public class Calc {

public void add(int summand1, int summand2) {
System.out.println(summand1 + summand2);
}

public static void main(String[] args) {
int summandA = 1;
int summandB = 2;
Calc calculator = new Calc();
System.out.print("1 + 2 = ");
calculator.add(summandA, summandB); 
// prints: 3
}

}
\end{lstlisting}
\end{frame}

\subsection{Return Value}
\begin{frame}[fragile]{Methods with Return Value}
A method without a return value is indicated by \textbf{void}:
\begin{lstlisting}
public void add(int summand1, int summand2) {
System.out.println(summand1 + summand2);
}
\end{lstlisting}
A method with an \textbf{int} as return value:
\begin{lstlisting}
public int add(int summand1, int summand2) {
return summand1 + summand2;
}
\end{lstlisting}
% TODO explain return statement
\end{frame}

\begin{frame}[fragile]{Calling Methods with a return value}
\begin{lstlisting}
public class Calc {

public int add(int summand1, int summand2) {
return summand1 + summand2;
}

public static void main(String[] args) {
Calc calculator = new Calc();
int sum = calculator.add(3, 8);
System.out.print("3 + 8 = " + sum); 
// prints: 3 + 8 = 11
}

}
\end{lstlisting}
\end{frame}

\subsection{Constructor}

\begin{frame}[fragile]{Constructors}
\begin{lstlisting}
public class Calc {

private int summand1;
private int summand2;

public Calc() {
summand1 = 0;
summand2 = 0;
}

}
\end{lstlisting}
A constructor gets called upon creation of the object
\end{frame}

\begin{frame}[fragile]{Constructors with Arguments}
\begin{lstlisting}
public class Calc {

private int summand1;
private int summand2;

public Calc(int x, int y) {
summand1 = x;
summand2 = y;
}

}
\end{lstlisting}
\begin{lstlisting}
[...]
Calc myCalc = new Calc(7, 9);
\end{lstlisting}

A constructor can have arguments as well!
\end{frame}

\section{Conclusion}
\subsection{An Example}

\begin{frame}{Let's build a car}
Create a car with doors, wheels, gas, seats,... \\
\vspace{1em}
Focus on:
\begin{description}
	\item[ID] unique id for each car
	\item[Car] gas, speed
	\item[Doors] can open/close
	\item[Wheels] air pressure, size,...
	\item[Seats] free, quality
	\item[\dots]
\end{description}
Also implement refueling, open/close doors, decrease/increase air pressure,...
% 	The more you think about it, the more complex this program becomes.
% 	Focus on the relevant things.
%	Think about how the objects can be in relation, this will be discussed later
%	Show prepared classes in Java
\end{frame}


\begin{frame}[fragile]{Class Student}
\begin{lstlisting}
public static void main(String[] args) {
Student peter = new Student();
peter.changeName("Peter");
}
\end{lstlisting}
\end{frame}

\end{document}